\chapter{Limits and colimits}\label{ch:limits}

\begin{assumption}\label{ass:ch6}
Throughout this chapter $\cC$ has an initial object $0$ and a terminal object
$1$. In \cref{cex:headless-products,cex:mono-headless} we consider headless
cats in categories without initial objects; this is stated explicitly there.
\end{assumption}

\section{Extremal objects}\label{sec:extremal}

We have already shown (\cref{prop:legless-initial}) that the initial objects of
$\Fel(\cC)$ are precisely the legless cats, and that the initial object is
strict.

\begin{proposition}\label{prop:terminal-cat}
The caudal cat $\Caudal=\iota(1)$ is a terminal object of $\Fel(\cC)$. A cat is
terminal if and only if it is unipedal and its leg is a terminal object of $\cC$.
\end{proposition}

\begin{proof}
For every cat $K'$, $\Fel(\cC)(K',\Caudal)\iso\prod_{\alpha}\cC(L'_\alpha,1)$
by \cref{cor:Fel-hom}, which is a singleton. Terminal objects are isomorphic
to $\Caudal$, and by \cref{prop:iso} the cats isomorphic to $\Caudal$ are
precisely the unipedal cats whose leg is isomorphic to $1$.
\end{proof}

\begin{lemma}\label{lem:legs-representable}
The leg functor is representable by the cephalic cat: there is a natural
isomorphism $\legs\iso\Fel(\cC)(\Cephalic,-)$, whose component at $K$ sends
$\alpha\in\legs(K)$ to the morphism $(\alpha,{!_{L_\alpha}})\colon\Cephalic\to K$.
\end{lemma}

\begin{proof}
By \cref{cor:Fel-hom}, $\Fel(\cC)(\Cephalic,K)\iso\coprod_{\alpha}\cC(0,L_\alpha)$,
and each $\cC(0,L_\alpha)$ is a singleton. Naturality: for $(\varphi,F)\colon K\to K'$,
$(\varphi,F)\circ(\alpha,!)=(\varphi(\alpha),F_\alpha\circ{!})=(\varphi(\alpha),!)$.
\end{proof}

\begin{remark}\label{rem:head-detects}
The head of any particular cat carries no information. What
\cref{lem:legs-representable} uses is the existence of an initial object in
$\cC$, which makes the set of legs of a cat \emph{representable}, so that the
leg functor preserves all limits. In \cref{prop:head-represent} we show that,
for headless cats in a category $\cC$ with a terminal object, the leg functor is
representable if and only if $\cC$ has an initial object.
\end{remark}

\section{Coproducts and extensivity}\label{sec:coproducts}

\begin{proposition}\label{prop:coproducts}
Let $(K_i)_{i\in I}$ be a small family of cats, where $K_i$ has legs
$(L_{i,\alpha})_{\alpha\in A_i}$. The cat
\[
\coprod_{i\in I}K_i\defeq\bigl\langle L_{i,\alpha}\bigr\rangle_{(i,\alpha)\in\coprod_iA_i},
\]
with coprojections $\inj_i=(\alpha\mapsto(i,\alpha),\id)\colon K_i\to\coprod_iK_i$,
is a coproduct of the family. In particular $\Fel(\cC)$ has all small
coproducts, and $\legs$ preserves them.
\end{proposition}

\begin{proof}
By \cref{cor:Fel-hom}, for any cat $M$ with legs $(M_\beta)_{\beta\in B}$,
\begin{align*}
\Fel(\cC)\Bigl(\coprod_iK_i,M\Bigr)&\iso\prod_{(i,\alpha)}\coprod_\beta\cC(L_{i,\alpha},M_\beta)
=\prod_i\prod_{\alpha\in A_i}\coprod_\beta\cC(L_{i,\alpha},M_\beta)\\
&\iso\prod_i\Fel(\cC)(K_i,M),
\end{align*}
and the composite bijection is precomposition with the coprojections.
\end{proof}

\begin{corollary}\label{cor:leg-decomposition-preview}
Every cat $K$ is the coproduct of its unipedal subcats: the leg inclusions
$j_\alpha\colon\iota(L_\alpha)\to K$ exhibit $K$ as $\coprod_{\alpha}\iota(L_\alpha)$.
\end{corollary}

\begin{proof}
Apply \cref{prop:coproducts} to the family $(\iota L_\alpha)_\alpha$; the
resulting coproduct is isomorphic to $K$ by an isomorphism with bijective gait
$(\alpha,*)\mapsto\alpha$ and identity components, under which the
coprojections correspond to the $j_\alpha$.
\end{proof}

Recall that a category $\cE$ with small coproducts is (infinitary)
\emph{extensive} if, for every small family $(X_i)_{i\in I}$ of objects, the
functor
\begin{equation}\label{eq:extensive}
\coprod\colon\prod_{i\in I}\cE/X_i\too\cE\Big/\coprod_iX_i,\qquad (Y_i\to X_i)_i\longmapsto\Bigl(\coprod_iY_i\to\coprod_iX_i\Bigr)
\end{equation}
is an equivalence of categories. For the finitary notion and its theory see
\cite{CarboniLackWalters1993}.

\begin{proposition}\label{prop:extensive}
$\Fel(\cC)$ is extensive. Moreover coproducts in $\Fel(\cC)$ are disjoint: for
$i\neq j$ the pullback of $\inj_i$ and $\inj_j$ exists and is the legless cat.
\end{proposition}

\begin{proof}
Let $(K_i)_{i\in I}$ be a family with leg sets $A_i$, and let
$(\gamma,G)\colon M\to\coprod_iK_i$ be an object of the slice over the
coproduct, $M$ having leg set $B$. For each $i$ let $B_i=\gamma^{-1}(\{i\}\times A_i)$
and let $M_i$ be the cat $\langle M_\beta\rangle_{\beta\in B_i}$ with the morphism
$G_i\colon M_i\to K_i$ given by $\beta\mapsto\gamma(\beta)$ (read in $A_i$) and
the components $G_\beta$. Since $B=\coprod_iB_i$, there is an isomorphism over
$\coprod_iK_i$ between $\coprod_i(M_i\to K_i)$ and $(M\to\coprod_iK_i)$ with
bijective gait and identity components; so \eqref{eq:extensive} is essentially
surjective. A morphism over $\coprod_iK_i$ between $\coprod_iM_i$ and
$\coprod_iM'_i$ has a gait compatible with the maps to $\coprod_iA_i$, hence
maps $B_i$ into $B'_i$ for each $i$; it is therefore the coproduct of a unique
family of morphisms $M_i\to M'_i$ over $K_i$. So \eqref{eq:extensive} is fully
faithful. For disjointness, a cone $(P\to K_i,P\to K_j)$ over the cospan
$(\inj_i,\inj_j)$ must have a gait sending each leg of $P$ into both
$\{i\}\times A_i$ and $\{j\}\times A_j$, which are disjoint; hence $P$ is
legless, and the legless cat, being initial, is the pullback.
\end{proof}

\section{Limits}\label{sec:limits}

\begin{theorem}\label{thm:limits}
Let $\cJ$ be a small category. Then $\Fel(\cC)$ has all limits of shape $\cJ$
if and only if $\cC$ does. In that case, the limit of a diagram
$D\colon\cJ\to\Fel(\cC)$, $j\mapsto K_j=\langle L_{j,a}\rangle_{a\in A_j}$,
$(u\colon j\to j')\mapsto(\varphi_u,F_u)$, is
\[
\lim_{j}K_j=\Bigl\langle\,\lim_{j\in\cJ}L_{j,a_j}\Bigr\rangle_{(a_j)_j\in\lim_jA_j},
\]
where $\lim_jA_j$ is computed in $\Set$ and, for $a=(a_j)_j\in\lim_jA_j$, the
limit $\lim_jL_{j,a_j}$ is that of the diagram $D_a\colon\cJ\to\cC$ given by
$j\mapsto L_{j,a_j}$ and $u\mapsto(F_u)_{a_j}$. The leg functor preserves all
limits that exist.
\end{theorem}

\begin{proof}
The leg functor preserves all limits because it is representable
(\cref{lem:legs-representable}).

\emph{Sufficiency.} Suppose $\cC$ has limits of shape $\cJ$. For
$a=(a_j)\in\lim_jA_j$ we have $\varphi_u(a_j)=a_{j'}$ for every $u\colon j\to j'$,
so $(F_u)_{a_j}\colon L_{j,a_j}\to L_{j',a_{j'}}$ is well typed, and
functoriality of $D$ (that is, $(\varphi_{vu},F_{vu})=(\varphi_v,F_v)(\varphi_u,F_u)$,
whose components give $(F_{vu})_{a_j}=(F_v)_{a_{j'}}(F_u)_{a_j}$) shows that $D_a$
is a functor. Choose a limit $(\Lambda_a,(p_{a,j})_j)$ of $D_a$ and set
$\Lambda=\langle\Lambda_a\rangle_{a\in\lim A_j}$. Define $P_j\colon\Lambda\to K_j$
to have gait the projection $\pr_j\colon\lim_jA_j\to A_j$ and components
$p_{a,j}$. These form a cone: $(\varphi_u,F_u)P_j$ has gait $\varphi_u\pr_j=\pr_{j'}$
and components $(F_u)_{a_j}p_{a,j}=p_{a,j'}$. Let
$(Q_j=(\psi_j,G_j)\colon M\to K_j)_j$ be another cone, $M=\langle M_b\rangle_{b\in B}$.
The gaits $\psi_j$ form a cone in $\Set$, so there is a unique
$\psi\colon B\to\lim_jA_j$ with $\pr_j\psi=\psi_j$. For each $b$, the
morphisms $(G_j)_b\colon M_b\to L_{j,\psi_j(b)}$ form a cone over
$D_{\psi(b)}$, by the cone condition on the $Q_j$; hence there is a unique
$g_b\colon M_b\to\Lambda_{\psi(b)}$ with $p_{\psi(b),j}g_b=(G_j)_b$. The morphism
$(\psi,(g_b))\colon M\to\Lambda$ satisfies $P_j(\psi,g)=Q_j$, and it is the
unique such morphism because its gait and components are forced.

\emph{Necessity.} Suppose $\Fel(\cC)$ has limits of shape $\cJ$ and let
$E\colon\cJ\to\cC$ be a diagram. Let $\Lambda$ be a limit of $\iota E$. Since
$\legs$ preserves limits, $\legs(\Lambda)\iso\lim_j\legs(\iota E_j)=\lim_j1$,
which is a singleton. Hence $\Lambda$ is unipedal, so $\Lambda\iso\iota(X)$ for
some $X$ (\cref{prop:iota-ff}), and the limit cone lies, up to isomorphism, in
the image of the fully faithful functor $\iota$. Fully faithful functors
reflect limits, so $X$ is a limit of $E$.
\end{proof}

\begin{corollary}\label{cor:limit-formulas}
Suppose the relevant limits exist in $\cC$.
\begin{enumerate}
\item The product of cats $K=\langle L_\alpha\rangle_A$ and $K'=\langle L'_\beta\rangle_B$ is
  $K\times K'=\langle L_\alpha\times L'_\beta\rangle_{(\alpha,\beta)\in A\times B}$.
\item The equalizer of $(\varphi,F),(\psi,G)\colon K\to K'$ is
  $\langle\mathrm{Eq}(F_\alpha,G_\alpha)\rangle_{\alpha\in A_0}$ where
  $A_0=\{\alpha:\varphi(\alpha)=\psi(\alpha)\}$.
\item The pullback of $K\to M\leftarrow K'$ with gaits $\varphi\colon A\to\legs M$
  and $\psi\colon B\to\legs M$ has legs $L_\alpha\times_{M_{\varphi(\alpha)}}L'_\beta$
  indexed by $\{(\alpha,\beta)\in A\times B:\varphi(\alpha)=\psi(\beta)\}$.
\end{enumerate}
\end{corollary}

\begin{proof}
These are the cases of \cref{thm:limits} for the discrete category with two
objects, the parallel pair, and the cospan.
\end{proof}

\begin{corollary}\label{cor:kit-limits}
The full subcategory $\Kit(\cC)$ of kittens is closed under all finite limits
that exist in $\Fel(\cC)$.
\end{corollary}

\begin{proof}
By \cref{thm:limits}, the leg set of a limit of a finite diagram of kittens is
the limit of a finite diagram of finite sets, which is a subset of a finite
product of finite sets.
\end{proof}

\begin{remark}\label{rem:chloe-limits}
The full subcategory of cats satisfying the Chloe Condition (\cref{def:chloe})
is closed under coproducts (\cref{prop:chloe-basic}) but not, in general, under
products. In the category $\Graph$ of \cref{ex:connected}, which has an initial
object, a terminal object (a vertex with a loop) and finite products, the cat
$\iota(K_2)$, whose leg is the graph with two vertices joined by an edge, is
Chloe, while $\iota(K_2)\times\iota(K_2)\iso\iota(K_2\times K_2)$ is not, because
$K_2\times K_2$ is disconnected (\cref{ex:gideon-tensor}). In $\Set$ and $\Top$,
where products of connected objects are connected and the terminal object is
connected, Chloe cats are closed under finite products.
\end{remark}


\begin{remark}\label{rem:limits-honest}
The sufficiency half of \cref{thm:limits} is the elementary computation of
limits in $\Fam(\cC)$. The necessity half depends on the initial object of
$\cC$, through \cref{lem:legs-representable}. The following counterexample
shows that it fails for headless cats.
\end{remark}

\begin{counterexample}\label{cex:headless-products}
Let $\cC$ be the poset $\{x,y,\top\}$ with $x<\top$ and $y<\top$ and $x,y$
incomparable. It has a terminal object $\top$ and no initial object, and the
product $x\times y$ (that is, the meet $x\wedge y$) does not exist, since $x$
and $y$ have no common lower bound. Nevertheless the category of headless
cats $\Fel^{\hl}(\cC)\eqv\Fam(\cC)$ has all binary products. For families
$(X_a)_{a\in A}$ and $(Y_b)_{b\in B}$, let
$P=\{(a,b)\in A\times B: X_a\wedge Y_b\text{ exists}\}$, which excludes exactly
the pairs with $\{X_a,Y_b\}=\{x,y\}$, and set $Z_{(a,b)}=X_a\wedge Y_b$, with
projections $(P,Z)\to(A,X)$ and $(P,Z)\to(B,Y)$ given by $(a,b)\mapsto a$ and
$(a,b)\mapsto b$ and the inequalities $Z_{(a,b)}\le X_a$, $Z_{(a,b)}\le Y_b$. For an
object $W$ of $\cC$,
\begin{align*}
\Fam(\cC)\bigl((1,W),(P,Z)\bigr)&=\{(a,b)\in P:W\le X_a\wedge Y_b\}\\
&=\{(a,b)\in A\times B: W\le X_a,\ W\le Y_b\},
\end{align*}
because $W\le X_a$ and $W\le Y_b$ force $X_a\wedge Y_b$ to exist. The right-hand
side is $\Fam(\cC)((1,W),(A,X))\times\Fam(\cC)((1,W),(B,Y))$, naturally in $W$,
and the bijection is composition with the two projections.
Since every family is a coproduct of one-element families and
$\Fam(\cC)(\coprod_i(1,W_i),-)\iso\prod_i\Fam(\cC)((1,W_i),-)$, it follows that
$(P,Z)$ is a product. In particular the product of $(1,x)$ and $(1,y)$ is the
empty family.
\end{counterexample}

\section{Colimits}\label{sec:colimits}

Recall that a category is \emph{connected} if it is nonempty and any two
objects are joined by a zigzag, and that a \emph{connected colimit} is a
colimit of a diagram indexed by a small connected category.

\begin{theorem}\label{thm:colimits}
The category $\Fel(\cC)$ has all small colimits if and only if $\cC$ has all
small connected colimits. In that case, the colimit of a diagram
$D\colon\cJ\to\Fel(\cC)$, with notation as in \cref{thm:limits}, is computed
as follows. Let $\cG=\El(\legs D)$ be the category of elements of the functor
$\legs D\colon\cJ\to\Set$, with objects $(j,a)$, $a\in A_j$, and morphisms
$u\colon(j,a)\to(j',\varphi_u(a))$. For each connected component $c$ of $\cG$
let $D_c\colon c\to\cC$ be the functor $(j,a)\mapsto L_{j,a}$,
$u\mapsto(F_u)_a$. Then
\[
\colim_jK_j=\Bigl\langle\colim D_c\Bigr\rangle_{c\in\pi_0(\cG)},
\]
and $\legs$ preserves all colimits that exist.
\end{theorem}

\begin{proof}
The leg functor preserves all colimits because it is a left adjoint
(\cref{cor:first-string}).

\emph{Sufficiency.} Suppose $\cC$ has small connected colimits. Each component
$c$ of $\cG$ is a small connected category, so $\colim D_c$ exists; let
$q_{(j,a)}\colon L_{j,a}\to\colim D_c$ be its coprojections, for $(j,a)\in c$.
By \cref{lem:colim-set}, $\colim_jA_j\iso\pi_0(\cG)$ with coprojections
$a\mapsto[(j,a)]$. Let $\Xi=\langle\colim D_c\rangle_{c\in\pi_0(\cG)}$ and define
$Q_j\colon K_j\to\Xi$ with gait $a\mapsto[(j,a)]$ and components $q_{(j,a)}$.
These form a cocone, because $q_{(j',\varphi_u(a))}\circ(F_u)_a=q_{(j,a)}$ is the
cocone condition for $D_c$. Given another cocone
$(R_j=(\rho_j,G_j)\colon K_j\to M)_j$ with $M=\langle M_b\rangle_{b\in B}$, the
gaits form a cocone in $\Set$ and induce a unique $\rho\colon\pi_0(\cG)\to B$
with $\rho([(j,a)])=\rho_j(a)$. For each component $c$, the morphisms
$(G_j)_a\colon L_{j,a}\to M_{\rho(c)}$, for $(j,a)\in c$, form a cocone on
$D_c$ by the cocone condition on the $R_j$, and so induce a unique
$g_c\colon\colim D_c\to M_{\rho(c)}$. The morphism $(\rho,(g_c))\colon\Xi\to M$
is the unique morphism compatible with the cocones.

\emph{Necessity.} Suppose $\Fel(\cC)$ has small colimits, and let
$E\colon\cJ\to\cC$ with $\cJ$ small and connected. Let $\Xi$ be a colimit of
$\iota E$. Since $\legs$ preserves colimits, $\legs(\Xi)\iso\colim_\cJ1\iso\pi_0(\cJ)$,
a singleton by \cref{lem:colim-set}. Hence $\Xi\iso\iota X$ for some $X$, and
$X$ is a colimit of $E$ because $\iota$ is fully faithful.
\end{proof}

\begin{remark}\label{rem:connected-colimits}
The same proof shows that $\Fel(\cC)$ has all small connected colimits if and
only if $\cC$ does; since $\Fel(\cC)$ always has small coproducts and every
colimit is a coproduct of connected colimits (over the components of the
indexing category), this is equivalent to \cref{thm:colimits}. The statement
for colimits of a \emph{fixed} connected shape is more delicate, because the
components of $\cG$ need not have the same shape as $\cJ$, nor be finite when
$\cJ$ is. The following counterexample shows that finite cocompleteness of
$\cC$ does not imply that $\Fel(\cC)$ has coequalizers.
\end{remark}

\begin{counterexample}\label{cex:coequalizer-rationals}
Let $\cC$ be the totally ordered set $\mathbb Q\cap[0,1]$, regarded as a
category. It has an initial object $0$, a terminal object $1$, and all finite
colimits (finite joins are maxima, and coequalizers are trivial in a poset).
Let $(q_n)_{n\in\mathbb N}$ be a strictly increasing sequence in $\cC$
converging to the irrational number $r=1/\sqrt2$, and let
$K=\langle q_n\rangle_{n\in\mathbb N}$. The morphisms
\[
f=(\id_{\mathbb N},(q_n\le q_n)_n),\qquad g=(n\mapsto n+1,(q_n\le q_{n+1})_n)\colon K\too K
\]
have no coequalizer in $\Fel(\cC)$. Indeed, suppose $e\colon K\to Q$ were a
coequalizer. Then $\legs(Q)$ is the coequalizer in $\Set$ of $\id$ and the
successor map on $\mathbb N$, which is a singleton; so $Q\iso\iota(z)$ for some
$z\in\cC$, and $e$ amounts to the inequalities $q_n\le z$. For every
$w\in\cC$, the morphisms $K\to\iota(w)$ coequalizing $f$ and $g$ are all
morphisms $K\to\iota(w)$, which exist if and only if $q_n\le w$ for all $n$,
that is, if and only if $w\geq r$. The universal property therefore requires
$z\le w$ for all rational $w\ge r$ in $\cC$ and $z\ge q_n$ for all $n$, so
$z=r$, which is not in $\cC$.
\end{counterexample}

\begin{proposition}\label{prop:kit-colimits}
The category $\Kit(\cC)$ of kittens has all finite colimits if and only if
$\cC$ has all finite connected colimits; in that case they are computed as in
\cref{thm:colimits} and agree with the colimits in $\Fel(\cC)$.
\end{proposition}

\begin{proof}
If $\cJ$ is finite and each $K_j$ is a kitten, the category $\cG$ of
\cref{thm:colimits} is finite, so its components are finite connected
categories, and the construction in the proof of \cref{thm:colimits} uses
only finite connected colimits in $\cC$. The result is a kitten since
$\pi_0(\cG)$ is finite, and it is a colimit in $\Fel(\cC)$, hence in the full
subcategory $\Kit(\cC)$. Conversely, the adjunction $\legs\dashv\Delta_1$
restricts to finite sets and kittens, so $\legs\colon\Kit(\cC)\to\mathbf{FinSet}$
preserves finite colimits, and the necessity argument of \cref{thm:colimits}
applies verbatim to finite connected $\cJ$.
\end{proof}

\section{Monomorphisms and epimorphisms}\label{sec:mono-epi}

\begin{proposition}\label{prop:mono}
A feline morphism $(\varphi,F)\colon K\to K'$ is a monomorphism if and only if
$\varphi$ is injective and each $F_\alpha$ is a monomorphism in $\cC$.
\end{proposition}

\begin{proof}
Sufficiency: if $(\varphi,F)(\chi,G)=(\varphi,F)(\chi',G')$ then
$\varphi\chi=\varphi\chi'$, so $\chi=\chi'$, and $F_{\chi(b)}G_b=F_{\chi(b)}G'_b$,
so $G_b=G'_b$. Necessity: $\legs\iso\Fel(\cC)(\Cephalic,-)$ is representable,
hence preserves monomorphisms, so $\varphi$ is injective. If
$g,g'\colon X\to L_\alpha$ satisfy $F_\alpha g=F_\alpha g'$, then the morphisms
$j_\alpha\iota(g),j_\alpha\iota(g')\colon\iota X\to K$ have equal composites with
$(\varphi,F)$, namely $(\varphi(\alpha),F_\alpha g)$; hence they are equal, and
$g=g'$.
\end{proof}

\begin{counterexample}\label{cex:mono-headless}
The necessity of injectivity in \cref{prop:mono} depends on the head. Let
$\cC$ be the poset $\{x_1,x_2,z\}$ with $x_1<z$, $x_2<z$ and $x_1,x_2$
incomparable; it has a terminal object $z$ but no initial object. In
$\Fel^{\hl}(\cC)\eqv\Fam(\cC)$, consider the morphism
$m\colon(2,(x_1,x_2))\to(1,z)$ with the unique gait $2\to1$. For any family
$(B,W)$, a morphism $(B,W)\to(2,(x_1,x_2))$ must send $b$ to the index $i$ with
$W_b=x_i$, and exists only if no $W_b$ equals $z$; so there is at most one
morphism from any object to $(2,(x_1,x_2))$, and $m$ is a monomorphism. Its
gait is not injective.
\end{counterexample}

\begin{proposition}\label{prop:epi}
A feline morphism $(\varphi,F)\colon K\to K'$ is an epimorphism if and only if
$\varphi$ is surjective and, for each $\beta\in B$, the family
$(F_\alpha\colon L_\alpha\to L'_\beta)_{\alpha\in\varphi^{-1}(\beta)}$ is jointly
epimorphic in $\cC$ (that is, $gF_\alpha=g'F_\alpha$ for all
$\alpha\in\varphi^{-1}(\beta)$ implies $g=g'$).
\end{proposition}

\begin{proof}
Sufficiency: if $(\chi,G)(\varphi,F)=(\chi',G')(\varphi,F)$ then
$\chi\varphi=\chi'\varphi$, so $\chi=\chi'$ as $\varphi$ is surjective; and
$G_\beta F_\alpha=G'_\beta F_\alpha$ for all $\alpha\in\varphi^{-1}(\beta)$, so
$G_\beta=G'_\beta$. Necessity: suppose $\beta\notin\im\varphi$. Let $K''$ be
$K'\sqcup\iota(L'_\beta)$, and consider the coprojection
$\inj_1\colon K'\to K''$ and the morphism $\inj'\colon K'\to K''$ which agrees
with $\inj_1$ except that it sends the leg $\beta$ to the new leg, with
identity component. Then $\inj_1\neq\inj'$ but
$\inj_1(\varphi,F)=\inj'(\varphi,F)$. Hence $\varphi$ is surjective. Now fix
$\beta$ and $g,g'\colon L'_\beta\to Y$ with $gF_\alpha=g'F_\alpha$ for all
$\alpha\in\varphi^{-1}(\beta)$. Let $K''$ be obtained from $K'$ by replacing the
leg $\beta$ by $Y$, and let $G,G'\colon K'\to K''$ have identity gait, components
$g$ resp.\ $g'$ at $\beta$, and identity components elsewhere. Then
$G(\varphi,F)=G'(\varphi,F)$, so $G=G'$ and $g=g'$.
\end{proof}

\begin{remark}\label{rem:mono-epi-asymmetry}
The proofs of \cref{prop:mono,prop:epi} are not dual to one another. The
characterization of monomorphisms uses the head, through the
representability of $\legs$, and fails without it
(\cref{cex:mono-headless}). The characterization of epimorphisms uses only
test morphisms built from coproducts and from the replacement of one leg by
another object, and in particular does not use the tail. This
asymmetry reflects the fact that $\Fel(\cC)$ is a free completion under
coproducts, and not under products.
\end{remark}

\section{Preservation properties}\label{sec:preservation}

\begin{proposition}\label{prop:preservation}
\begin{enumerate}
\item The leg functor $\legs$ preserves all limits and colimits.
\item The one-legged inclusion $\iota$ preserves all limits that exist in $\cC$,
  and all connected colimits that exist in $\cC$; it preserves neither the
  initial object nor binary coproducts.
\item $\Delta_0$ preserves all colimits, and $\Delta_1$ preserves all limits and
  all colimits.
\end{enumerate}
\end{proposition}

\begin{proof}
(i) was shown in \cref{thm:limits,thm:colimits}. (ii) If $X$ is a limit of
$E\colon\cJ\to\cC$, the cone $\iota X\to\iota E$ is a limit cone. Indeed, the
sufficiency argument in the proof of \cref{thm:limits} uses only the limits of
the diagrams $D_a$ for $a\in\lim_jA_j$; for $D=\iota E$ we have $\lim_j1=1$ and
$D_*=E$, whose limit exists. Explicitly, a cone $M\to\iota E$ consists of the
unique gait $\legs M\to1$ and, for each leg of $M$, a cone over $E$, which
factors uniquely through $X$. Similarly, the sufficiency argument of
\cref{thm:colimits} with $\El(\legs\iota E)\iso\cJ$ connected uses only the
colimit of $E$, so $\iota$ preserves connected colimits that exist. The initial object
$0$ is sent to $\Cephalic$, which is not initial (\cref{warn:cephalic}); and
$\iota(X\sqcup Y)$ is unipedal while $\iota X\sqcup\iota Y$ is two-legged.
(iii) $\Delta_0$ is a left adjoint; $\Delta_1$ is a right adjoint
(\cref{cor:first-string}) and also a left adjoint (\cref{thm:adjoint-string}).
\end{proof}
